\documentclass{beamer}

\usepackage{amsmath}
\usepackage{amssymb}
\usepackage{graphicx}
\usepackage{subcaption}
%\usepackage{subfigure}
\usepackage{float}
\usepackage{multirow}
\usepackage{alphalph}
\usepackage{hyperref}
%\usepackage{lipsum}
%\usepackage{verbatim}

\usepackage{appendixnumberbeamer}
\usetheme{Frankfurt}
%\xdefinecolor{maroon}{cmyk}{0.15,1.00,0.39,0.69}
%\xdefinecolor{maroon}{cmyk}{0.27,0.86,0.60,0.30}
%\xdefinecolor{maroon}{cmyk}{0.000,1.00,1.00,0.498}
\usecolortheme[RGB={80,0,0}]{structure}

% Shortcut commands
% My Shortcut Commands
\newcommand{\fig}[1]{Fig.~\ref{#1}}                      % figure
\newcommand{\tbl}[1]{Table~\ref{#1}}                     % table

\newcommand{\be}{\begin{equation*}}   % numbered equation
\newcommand{\ee}{\end{equation*}}

\newcommand{\benum}{\begin{equation}}   % numbered equation
\newcommand{\eenum}{\end{equation}}

\newcommand{\bea}{\begin{eqnarray*}}  % numbered equation array
\newcommand{\eea}{\end{eqnarray*}}

\newcommand{\beanum}{\begin{eqnarray}}  % numbered equation array
\newcommand{\eeanum}{\end{eqnarray}}

\newcommand{\eqt}[1]{Eq. (\ref{#1})}  % Reference to one equation
\newcommand{\eqts}[1]{Eqs. (\ref{#1})}  % Reference to multiple equations 

\newcommand{\vect}[1]{\ensuremath{ \vec{\mathbf #1}}}  % bold faced with vector arrow above
\newcommand{\B}[1]{\ensuremath{B_{#1} }}			% B with a sub scripted argument

\newcommand{\p}{\ensuremath{ \partial}}			% shortcut partial derivative symbol

\newcommand{\abs}[1]{\ensuremath{\left\lvert #1 \right\rvert}}  % absolute value of argument (variable bar size)
\newcommand{\norm}[1]{\ensuremath{\left\lVert #1 \right\rVert}}  % norm of argument, varaible size

\newcommand{\omg}{\ensuremath{\vec{\Omega}}}
\newcommand{\BCSZ}{\ensuremath{\widetilde{\psi}_{BCSZ}}}
\newcommand{\BCSZH}{\ensuremath{\widehat{\psi}_{BCSZ}}}

% Equation Punctuation
\newcommand{\pec}{\, ,}
\newcommand{\pep}{\, .}
\setbeamertemplate{footline}{
\leavevmode%
\hbox{\hspace*{-0.06cm}
\begin{beamercolorbox}[wd=.45\paperwidth,ht=2.25ex,dp=1ex,center]{author in head/foot}%
	\usebeamerfont{author in head/foot}Maginot%~~(\insertshortinstitute)
\end{beamercolorbox}%
\begin{beamercolorbox}[wd=.3\paperwidth,ht=2.25ex,dp=1ex,center]{section in head/foot}%
PhD Defense %\usebeamerfont{section in head/foot}\insertshorttitle
\end{beamercolorbox}%
\begin{beamercolorbox}[wd=.25\paperwidth,ht=2.25ex,dp=1ex,right]{section in head/foot}%
	\usebeamerfont{section in head/foot}April 2015\hspace*{2em}
	\insertframenumber{} / \inserttotalframenumber\hspace*{2ex}
\end{beamercolorbox}}%
\vskip0pt%
}
\beamertemplatenavigationsymbolsempty
\beamertemplatetransparentcovered

\setbeamertemplate{bibliography item}[text]

\title{PhD Defense \\ Higher Order Radiative Transfer}
\author{Peter Maginot}\institute{Texas A\&M University- Department of Nuclear Engineering}
\date{April 2015}


\newif\ifplacelogo % create a new conditional
\placelogotrue % set it to true
%\logo{\ifplacelogo\includegraphics[height=1cm]{defense_presentation/CSGF_vert.pdf} \hspace{3.5in} \includegraphics[height=1cm]{defense_presentation/cert-logo.jpg}\fi}

\begin{document}
%---------------------------------------------------------------------------------------------
\placelogofalse
\begin{frame}
\titlepage


\end{frame}

%---------------------------------------------------------------------------------------------

\placelogofalse
%---------------------------------------------------------------------------------------------

%\begin{frame}
%\frametitle{Road to High Order Radiative Transfer}
%\tableofcontents[hideallsubsections]
%%\begin{enumerate}
%%\item Spatially Varying Cross Section Problems
%%\item Low Order DSA for High Order DFEM Transport
%%\item Strictly Non-Negative Bilinear Scheme for Unstrutured Quads
%%\end{enumerate}
%\end{frame}

\section{Goal}
\subsection{Objectives}
\begin{frame}
\frametitle{Holder}
\end{frame}

\section{Old Results}
\subsection{Self-Lumping DFEM}
\begin{frame}
\frametitle{Holder}
\end{frame}

\subsection{Spatially Varying Cross Section}
\begin{frame}
\frametitle{Holder}
\end{frame}

\section{Derivation}
\subsection{Temperature Linearization}
\begin{frame}
\frametitle{Holder}
\end{frame}

\subsection{MIP DSA for Radiative Transfer}
\begin{frame}
\frametitle{Holder}
\end{frame}

\section{Computer Science}
\subsection{Code Features}
\begin{frame}
\frametitle{Holder}
\end{frame}

\subsection{C++ Things}
\begin{frame}
\frametitle{Holder}
\end{frame}

\section{Results}
\subsection{MMS}
\begin{frame}
\frametitle{Holder}
\end{frame}

\subsection{Marhsak Wave}
\begin{frame}
\frametitle{Holder}
\end{frame}

\section{Conclusions}
\subsection{Near Term Future Work}
\begin{frame}
\frametitle{Holder}
\end{frame}

\subsection{Long Term Future Work}
\begin{frame}
\frametitle{Holder}
\end{frame}

\subsection{Acknowledgments}
\begin{frame}
\frametitle{Acknowledgments}
Thanks for your time!
\vspace{0.3in}
Portions of this work were funded by the Department of Energy CSGF program, administered by the Krell Institute, under grant DE-FG02-97ER25308.
\\
\vspace{0.3in}
Additional support was provided by the Department of Energy, National Nuclear Security Administration, under Award Number(s) DE-NA0002376.

\end{frame}

%%%%%%%%%%%%%%%%%%%%%%%%%%%%%%%%%%%%%%%%%%%%%%%%%%%%%%%%%%%%%%%%%%%%%%%%%%%%%%%%%

\end{document}
%------------------------------------------------------------------------------